\subsection{Inverse systems of measures}

DEFINITION 1. --- Let \(\mathcal{T} = (\mathrm{T}_i, p_{ij})\) be an inverse system of topological spaces indexed by I. One calls inverse system (resp. sub-inverse system)

of measures on \(\mathcal{T}\) a family \((\mu_i)_{i\in I}\) , where \(\mu_i\) is a bounded measure on \(\mathrm{T}_i\) for all \(i\in \mathrm{I}\) , and where \(\mu_i = p_{ij}(\mu_j)\) (resp. \(\mu_i\geqslant p_{ij}(\mu_j)\) ) for \(i\leqslant j\) .

PROPOSITION 3. --- Let there be given an inverse system of topological spaces \(\mathcal{T} = (\mathrm{T}_i, p_{ij})\) indexed by I, a topological space T, a coherent and separating family of continuous mappings \(p_i: \mathrm{T} \to \mathrm{T}_i\) (for \(i \in \mathrm{I}\) ) and a sub-inverse system \((\mu_i)_{i \in \mathrm{I}}\) of measures on \(\mathcal{T}\) . For every compact subset K of T, set

\[
\mathrm{J} (\mathrm{K}) = \inf _ {i \in \mathrm{I}} \mu_ {i} ^ {\bullet} (p _ {i} (\mathrm{K})).\tag{1}
\]

Then there exists a bounded measure \(\pi\) on T, and only one, such that \(\pi^{\bullet}(\mathrm{K}) = \mathrm{J}(\mathrm{K})\) for every compact subset K of T. One has \(\mu_i \geqslant p_i(\pi)\) for all \(i \in I\) , and \(\pi\) is the largest measure on T satisfying this condition.

Let us first prove that \(\mathrm{J}(\mathrm{K})\) is the limit of \(\mu_i^\bullet (p_i(\mathrm{K}))\) with respect to the section filter \(\mathfrak{F}\) of the directed preordered set I: for this, it suffices (GT, IV, §5, No.~2, Th. 2) to show that \(\mu_i^\bullet (p_i(\mathrm{K})) \geqslant \mu_j^\bullet (p_j(\mathrm{K}))\) for \(i \leqslant j\) ; now, setting \(\mu_{ij}' = p_{ij}(\mu_j)\) , we have \(\mu_{ij}' \leqslant \mu_i\) and \(p_j(\mathrm{K}) \subset \overline{p}_{ij}^{-1}(p_i(\mathrm{K}))\) , whence

\[
\mu_ {j} ^ {\bullet} \big (p _ {j} (\mathrm{K}) \big) \leqslant \mu_ {j} ^ {\bullet} \big (\overline{p}_{ij}^{-1} (p _ {i} (\mathrm{K})) \big) = (\mu_ {i j} ^ {\prime}) ^ {\bullet} \big (p _ {i} (\mathrm{K}) \big) \leqslant \mu_ {i} ^ {\bullet} \big (p _ {i} (\mathrm{K}) \big).
\]

Let us now pass to the study of the properties of the function J:

\begin{enumerate}
\def\labelenumi{\arabic{enumi})}
\item
  It is clear that \(\mathrm{J}(\mathbf{K})\leqslant \mathrm{J}(\mathbf{L})\) when \(\mathbf{K}\subset \mathbf{L}\) .
\item
  Let \(\mathbf{K}\) and \(\mathbf{L}\) be two compact subsets of \(\mathbf{T}\) . For every \(i \in \mathbf{I}\) , we have \(p_i(\mathbf{K} \cup \mathbf{L}) = p_i(\mathbf{K}) \cup p_i(\mathbf{L})\) , whence
\end{enumerate}

\[
\mu_ {i} ^ {\bullet} \left(p _ {i} (\mathrm{K} \cup \mathrm{L})\right) \leqslant \mu_ {i} ^ {\bullet} \left(p _ {i} (\mathrm{K})\right) + \mu_ {i} ^ {\bullet} \left(p _ {i} (\mathrm{L})\right);
\]

passing to the limit with respect to the filter \(\mathfrak{F}\) , we obtain \(\mathrm{J}(\mathrm{K} \cup \mathrm{L}) \leqslant \mathrm{J}(\mathrm{K}) + \mathrm{J}(\mathrm{L})\) .

\begin{enumerate}
\def\labelenumi{\arabic{enumi})}
\setcounter{enumi}{2}
\tightlist
\item
  Suppose that the compact sets K and L are disjoint. By Prop. 2 of No.~1, there exists an \(i \in I\) such that \(p_j(K) \cap p_j(L) = \varnothing\) for \(j \geqslant i\) . For \(j \geqslant i\) we therefore have
\end{enumerate}

\[
\mu_ {j} ^ {\bullet} \left(p _ {j} (\mathrm{K} \cup \mathrm{L})\right) = \mu_ {j} ^ {\bullet} \left(p _ {j} (\mathrm{K})\right) + \mu_ {j} ^ {\bullet} \left(p _ {j} (\mathrm{L})\right),
\]

whence \(\mathrm{J}(\mathrm{K}\cup\mathrm{L})=\mathrm{J}(\mathrm{K})+\mathrm{J}(\mathrm{L})\) on passing to the limit with respect to the filter F.

\begin{enumerate}
\def\labelenumi{\arabic{enumi})}
\setcounter{enumi}{3}
\tightlist
\item
  Let \((\mathbf{K}_{\alpha})_{\alpha \in \mathbf{A}}\) be a decreasing directed family of compact subsets of T, with intersection K. By Prop. 1 of No.~1, \(p_i(\mathbf{K}) = \bigcap_{\alpha \in \mathbf{A}} p_i(\mathbf{K}_{\alpha})\) and
\end{enumerate}

therefore \(\mu_i^\bullet (p_i(\mathbf{K})) = \inf_{\alpha \in \mathbf{A}}\mu_i^\bullet (p_i(\mathbf{K}_\alpha))\) for all \(i\in I\) (§1, No.~6, Cor. of Prop. 5). From this, one deduces

\[
\begin{array}{r l} \mathrm{J(K)} & = \inf _ {i \in \mathrm{I}} \mu_ {i} ^ {\bullet} (p _ {i} (\mathrm{K})) = \inf _ {i \in \mathrm{I}} \inf _ {\alpha \in \mathrm{A}} \mu_ {i} ^ {\bullet} (p _ {i} (\mathrm{K} _ {\alpha})) \\ & = \inf _ {\alpha \in \mathrm{A}} \inf _ {i \in \mathrm{I}} \mu_ {i} ^ {\bullet} (p _ {i} (\mathrm{K} _ {\alpha})) = \inf _ {\alpha \in \mathrm{A}} \mathrm{J(K} _ {\alpha}). \end{array}
\]

\begin{enumerate}
\def\labelenumi{\arabic{enumi})}
\setcounter{enumi}{4}
\tightlist
\item
  Let us choose an \(i \in I\) and set \(c = \mu_i^\bullet(T_i)\) . Then \(c\) is finite and \(J(K) \leqslant \mu_i^\bullet(p_i(K)) \leqslant \mu_i^\bullet(T_i)\) , thus \(J(K) \leqslant c\) for every compact set \(K\) in \(T\) . The preceding properties permit applying Th. 1 of §3, No.~1; we conclude that there exists one and only one bounded measure \(\pi\) on \(T\) such that \(\pi^\bullet(K) = J(K)\) for every compact subset \(K\) of \(T\) . For every \(i \in I\) , let us denote by \(\nu_i\) the measure on \(T_i\) that is the image of \(\pi\) under \(p_i\) . Let \(i \in I\) , \(A\) a compact subset of \(T_i\) , and \(\mathfrak{L}\) the set of compact subsets of \(\overline{p}_i^1(A)\) . By Remark 3 of §1, No.~2, we have \(\pi^\bullet(\overline{p}_i^1(A)) = \sup_{K \in \mathfrak{L}} \pi^\bullet(K)\) ; moreover, \(\nu_i^\bullet(A) = \pi^\bullet(\overline{p}_i^1(A))\) and \(J(K) = \pi^\bullet(K)\) for \(K \in \mathfrak{L}\) , whence \(\nu_i^\bullet(A) = \sup_{K \in \mathfrak{L}} J(K)\) . For \(K \in \mathfrak{L}\) , we have \(p_i(K) \subset A\) , whence \(J(K) \leqslant \mu_i^\bullet(p_i(K)) \leqslant \mu_i^\bullet(A)\) and finally \(\nu_i^\bullet(A) \leqslant \mu_i^\bullet(A)\) . Since \(A\) is an arbitrary compact set in \(T_i\) , we conclude that \(\nu_i \leqslant \mu_i\) . The last assertion of the proposition is obvious.
\end{enumerate}

Q.E.D.

THEOREM 1 (Prokhorov). --- Let \(\mathcal{T} = (\mathrm{T}_{i}, p_{ij})\) be an inverse system of topological spaces indexed by I, T a topological space and \((p_{i})_{i \in I}\) a coherent and separating family of continuous mappings \(p_{i}: \mathrm{T} \to \mathrm{T}_{i}\) . Finally, let \((\mu_{i})_{i \in I}\) be an inverse system of measures on \(\mathcal{T}\) .

For there to exist a bounded measure \(\mu\) on \(\mathbf{T}\) such that \(p_i(\mu) = \mu_i\) for all \(i \in \mathbf{I}\) , it is necessary and sufficient that the following condition be satisfied:

\begin{enumerate}
\def\labelenumi{(\Alph{enumi})}
\setcounter{enumi}{15}
\tightlist
\item
  for every \(\varepsilon > 0\) , there exists a compact subset \(K\) of \(T\) such that \(\mu_i^\bullet(T_i - p_i(K)) \leqslant \varepsilon\) for all \(i \in I\) .
\end{enumerate}

When this is so, the measure \(\mu\) is uniquely determined and

\[
\mu^ {\bullet} (\mathbf {K}) = \inf _ {i} \mu_ {i} ^ {\bullet} (p _ {i} (\mathbf {K}))\tag{2}
\]

for every compact set \(\mathbf{K}\) in \(\mathbf{T}\) .

Let us first prove the uniqueness of \(\mu\) . Let \(\mu\) be a bounded measure on T such that \(p_i(\mu) = \mu_i\) for all \(i \in I\) . Let K be a compact subset of T; by Prop. 2 of No.~1, the set K is the intersection of the decreasing directed family \(\left(\overline{p_i^1}(p_i(K))\right)_{i \in I}\) of closed subsets of T. By the Cor. of Prop. 5 of §1, No.~6, we therefore have

\[
\mu^ {\bullet} (\mathrm{K}) = \inf _ {i \in \mathrm{I}} \mu^ {\bullet} \left(\bar {p} _ {i} ^ {- 1} (p _ {i} (\mathrm{K}))\right) = \inf _ {i \in \mathrm{I}} \mu_ {i} ^ {\bullet} \left(p _ {i} (\mathrm{K})\right),
\]

which establishes the formula (2). Since two measures that coincide on the set of compact sets are equal (§1, No.~2, Cor. of Prop. 2), it follows that \(\mu\) is unique.

By Prop. 3, there exists a bounded measure \(\pi\) on \(\mathbf{T}\) such that \(\pi^{\bullet}(\mathbf{K}) = \inf_{i\in I}\mu_i^\bullet (p_i(\mathbf{K}))\) for every compact subset \(\mathbf{K}\) of \(\mathbf{T}\) . By formula (2), the existence of a bounded measure \(\mu\) on \(\mathbf{T}\) such that \(p_i(\mu) = \mu_i\) for all \(i\in I\) is therefore equivalent to the relation:

\[
\left(\mathrm{P} ^ {\prime}\right) p _ {i} (\pi) = \mu_ {i} \text { for all } i \in \mathrm{I}.
\]

For \(i \leqslant j\) , we have \(\mu_i = p_{ij}(\mu_j)\) , whence \(\mu_i^\bullet (\mathbf{T}_i) = \mu_j^\bullet (\mathbf{T}_j)\) ; since I is directed, there exists a finite number \(c \geqslant 0\) such that \(\mu_i^\bullet (\mathbf{T}_i) = c\) for all \(i \in I\) . By Prop. 3, the measure \(\mu_i - p_i(\pi)\) is positive, hence is zero if and only if its total mass is zero, that is, if \(\mu_i(\mathbf{T}_i) = p_i(\pi)^\bullet (\mathbf{T}_i)\) . Since \(p_i(\pi)^\bullet (\mathbf{T}_i) = \pi^\bullet (\mathbf{T})\) , the condition \((\mathbf{P}')\) is thus equivalent to \(\pi^\bullet (\mathbf{T}) = c\) , that is (§1, No.~2, Remark 3) to the property:

\((\mathbf{P}^{\prime \prime})\sup_{K\in \mathfrak{K}}\pi^{\bullet}(\mathbf{K}) = c\) , where \(\mathfrak{K}\) is the set of compact subsets of T.

Now, for \(\mathbf{K} \in \mathfrak{K}\) , we have

\[
\pi^ {\bullet} (\mathrm{K}) = \inf _ {i \in \mathrm{I}} \mu_ {i} ^ {\bullet} (p _ {i} (\mathrm{K})) = c - \sup _ {i \in \mathrm{I}} \mu_ {i} ^ {\bullet} (\mathrm{T} _ {i} - p _ {i} (\mathrm{K}))
\]

and this formula immediately implies the equivalence of (P) and \((\mathbf{P}^{\prime \prime})\) .

Q.E.D.

Let \((\mathrm{T}_{i}, p_{ij})\) be an inverse system of topological spaces. Set \(\mathrm{T} = \varprojlim \mathrm{T}_{i}\) and denote by \(p_{i}\) the canonical mapping of \(\mathrm{T}\) into \(\mathrm{T}_{i}\) . Generalizing Def. 2 of Ch. III, §4, No.~5, we shall say that a bounded measure \(\mu\) on \(\mathrm{T}\) is the inverse limit of an inverse system \((\mu_{i})_{i \in \mathrm{I}}\) of measures if \(\mu_{i} = p_{i}(\mu)\) for all \(i \in \mathrm{I}\) . Th. 1 provides a criterion for the existence of inverse limits of measures. When the spaces \(\mathrm{T}_{i}\) are compact, and the mappings \(p_{ij}\) surjective, \(\mathrm{T}\) is compact and \(p_{i}(\mathrm{T}) = \mathrm{T}_{i}\) for every \(i \in \mathrm{I}\) ; the condition (P) is therefore fulfilled, and in this case we recover Prop. 8, (iv) of Ch. III, §4, No.~5.

Remark. --- Let \((\mu_{i})_{i\in\mathbb{I}}\) be an inverse system of measures on the inverse system of spaces \(\mathcal{T}=(\mathrm{T}_{i},p_{ij})\) . Assume given a topological space \(T'\) and continuous mappings \(p_{i}^{\prime}:T^{\prime}\to T_{i}\) ; assume that the family \((p_{i}^{\prime})_{i\in\mathbb{I}}\) is coherent, but not necessarily separating. If Prokhorov's condition (P) is satisfied by the family \((p_{i}^{\prime})_{i\in\mathbb{I}}\) , there exists a measure \(\mu^{\prime}\) (not necessarily unique) on \(T^{\prime}\) such that \(p_{i}^{\prime}(\mu^{\prime})=\mu_{i}\) for all \(i\in I\) .

For, set \(\mathbf{T} = \varprojlim \mathbf{T}_i\) and \(p' = (p_i')_{i\in I}\) , and denote by \(p_i\) the canonical mapping of \(\mathbf{T}\) into \(\mathbf{T}_i\) ; Prokhorov's condition is satisfied by \(\mathbf{T}\) and the \(p_i\) , because \(p_i(p'(K')) = p'_i(K')\) and \(p'(K')\) is compact in \(\mathbf{T}\) for every compact subset \(\mathbf{K}'\) of \(\mathbf{T}'\) . By Th. 1, there exists a bounded measure \(\mu\) on \(\mathbf{T}\) such that \(p_i(\mu) = \mu_i\) for all \(i\in I\) . Let \(\mathbf{K}'\) be a compact set in \(\mathbf{T}'\) ; then \(\mu^\bullet (p'(K')) = \inf_{i\in I}\mu_i^\bullet (p_i'(K'))\) ,

whence

\[
\mu^ {\bullet} \left(\mathrm{T} - p ^ {\prime} \left(\mathrm{K} ^ {\prime}\right)\right) = \sup _ {i \in \mathrm{I}} \mu_ {i} ^ {\bullet} \left(\mathrm{T} _ {i} - p _ {i} ^ {\prime} \left(\mathrm{K} ^ {\prime}\right)\right).
\]

Let \(\varepsilon > 0\) ; since Prokhorov's condition (P) is satisfied by the \(p_i'\) , one can therefore find a compact subset \(K'\) of \(T'\) such that \(\mu^\bullet(T - p'(K')) \leqslant \varepsilon\) . Prop. 8 of §2, No.~4 then establishes the existence of a bounded measure \(\mu'\) on \(T'\) with \(\mu = p'(\mu')\) , whence \(\mu_i = p_i(\mu) = p_i(p'(\mu')) = p_i'(\mu')\) for all \(i \in I\) .

